\documentclass[10pt,a4paper]{amsart}
\usepackage[margin=19mm]{geometry}
\usepackage{amsmath,amssymb,mathtools}
\usepackage[scr=rsfs,cal=euler]{mathalfa}
\usepackage{xfrac}
\usepackage[unicode,hidelinks,bookmarks=false]{hyperref}
\newcommand{\ie}{i.e.\ }
\newcommand{\cf}{cf.\ }
\newcommand{\resp}{respectively\ }
\newcommand{\Subsection}{\subsection}
\providecommand{\nobreakdash}{\nobreak\hbox{-}}
\setlength{\emergencystretch}{2em}
\multlinegap0pt
\newcommand{\A}{\mathbb{A}}
\usepackage{xcolor}
\usepackage{amssymb}
\usepackage[shortlabels]{enumitem}
\newtheorem{lem}{Lemma}
\newcommand{\N}{\mathbb{N}}
\newcommand{\R}{\mathbb{R}}
\DeclareMathOperator{\Spec}{Spec}
\newcommand{\gp}{\mathrm{gp}}
\newcommand{\ol}{\overline}
\title{$\A^1$-homotopy theory of log schemes}
\author{Doosung Park}
\date{Source: arXiv:2205.14750v4 (CC BY 4.0)}
\begin{document}
\maketitle

\noindent\emph{Source and licence.} $\A^1$-homotopy theory of log schemes. Version arXiv:2205.14750v4 (CC BY 4.0), distributed by arXiv under the Creative Commons Attribution 4.0 International licence, http://creativecommons.org/licenses/by/4.0/, as recorded in the arXiv licence metadata for this exact version and shown on the abstract page https://arxiv.org/abs/2205.14750v4. Full licence notice: \texttt{licenses/arxiv-2205.14750-CC-BY-4.0.txt}.\par\smallskip

\noindent\textit{Editorial scope.} Counted unit: the article's lem lemma.16 (source0.tex lines 3021--3029). The article's complete original proof of it is delivered with it (source0.tex lines 3030--3087, one \texttt{proof} environment of 58 lines). No mathematics is written, repaired or re-derived: the statement and the proof are the article's own lines, in the article's own notation, and no retained line is substituted or re-flowed.  No further source-local context is needed: every cross-reference of every retained line resolves inside the two delivered spans.  Nothing was omitted from any retained span, so the package carries no omission record.  No retained line contains a citation, so the wrapper carries no bibliography and no bibliography entry.  No retained line contains a figure environment, an image-inclusion command or drawing code, so no image is delivered and the package declares no asset dependency.  Editorial formatting only, in addition: an AMS \texttt{amsart} preamble that restates the article's own declarations and macros used by the retained lines, namely the environments \texttt{lem} on separate counters, together with the standard packages those lines need.  The retained environments print in this document as Lemma 1 in order of appearance, whereas the article prints the same items with the numbering of its own full document.  The complete native source of the article as distributed in the arXiv e-print for this exact version is bundled unchanged as \texttt{sources/126288/source0.tex}.
\begin{lem}
\label{lemma.16}
Let $\theta\colon \N\to P$ be a nontrivial map of fs monoids.
If $P^\gp$ is torsion free and $P\neq P^*$, then the topological space
\[
|\Spec(P)-\partial_{\Spec(\N)}\Spec(P)|_{\R}-\{0\}
\]
is acyclic.
\end{lem}

\begin{proof}
If $\theta$ is vertical, then
$|\Spec(P)-\partial_{\Spec(\N)}\Spec(P)|_{\R}-\{0\}=|\Spec(P)|_{\R}-\{0\}$ is contractible since $|\Spec(P)|_{\R}-\{0\}$ is convex.

Assume $\theta$ is not vertical.
Let $f\colon |\Spec(P)|_{\R}\to |\Spec(\N)|_{\R}\simeq \R_{\geq 0}$ be the morphism induced by $\theta$.
Choose any map $\eta\colon \N\to P$ such that $\eta(1)$ is not contained in any proper face of $P$.
This induces a map $h\colon |\Spec(P)|_{\R}\to \R_{\geq 0}$ such that $h^{-1}(0)=0$.
We set
\[
V:=f^{-1}(0)\cap h^{-1}(1)
\text{ and }
W:=|\Spec(P)-\partial_{\Spec(\N)}\Spec(P)|_{\R}\cap h^{-1}(1).
\]
For a cone $\sigma\in \Spec(P)$, we have $\sigma\in \Spec(P)-\partial_{\Spec(\N)}\Spec(P)$ if and only if $\sigma$ does not contain a ray $r$ with $f(r)=0$. 
Hence $f(x)>0$ for all $x\in |\Spec(P)-\partial_{\Spec(\N)}\Spec(P)|_{\R}-\{0\}$, so we have $V\cap W=\emptyset$.
There is a homeomorphism
\[
|\Spec(P)-\partial_{\Spec(\N)}\Spec(P)|_{\R}-\{0\}
\simeq
W\times \R_{>0}.
\]
Hence it suffices to show that $W$ is acyclic.

Let $X$ be the boundary of $h^{-1}(1)$.
Since $h^{-1}(1)$ is a convex polytope, $X$ is homeomorphic to $S^{n-2}$, where $n$ is the rank of $\ol{P}^\gp$.
We assumed that $\theta$ is not vertical, so we have $W\subset X$.
By the Alexander duality
\[
\tilde{H}_q(X-W)
\simeq
\tilde{H}^{n-q-3}(W)
\]
that holds for every integer $q\geq 0$, we only need to show that $U:=X-W$ is acyclic.

Let $\Gamma$ be the set of cones $\sigma$ of $\Spec(P)$ such that there exists two rays $r_1$ and $r_2$ of $\sigma$ with $f(r_1)=0$ and $f(r_2)\neq 0$.
Equivalently, $|\sigma|_{\R}\cap V,|\sigma|_{\R}\cap W\neq \emptyset$.
We set $u(\sigma):=|\sigma|_{\R}\cap U$.
Let $\{\sigma_1,\ldots,\sigma_m\}$ Let $I:=\{i_1,\ldots,i_r\}$ be a nonempty subset of $\{1,\ldots,m\}$, and we set $\sigma_I:=\sigma_{i_1}\cap \cdots \cap \sigma_{i_r}$.
If $u(\sigma_I)=\emptyset$, then $u(\sigma_I)\cap V=\emptyset$.
On the other hand, if $u(\sigma_I)\neq \emptyset$, then $\sigma_I$ contains a ray $r$ with $f(r)=0$.
Hence we have $u(\sigma_I)\cap V\neq \emptyset$.
Since $u(\sigma_I)$ and $u(\sigma_I)\cap V$ are convex, they are contractible.
Hence the inclusion
\[
u(\sigma_I)\cap V
\to
u(\sigma_I)
\]
is a homotopy equivalence.

Since 
$\{u(\sigma_1),\ldots,u(\sigma_m),V\}$ (resp.\ $\{u(\sigma_1)\cap V,\ldots,u(\sigma_m)\cap V,V\}$) is a closed cover of $U$ (resp.\ $V$), the homomorphism of singular cohomology groups $H_*(V)\to H_*(U)$ induced by the inclusion $V\to U$ is an isomorphism by 
Mayer-Vietoris.
The assumption that $\theta$ is nontrivial implies that $V$ is not empty.
Furthermore, $V$ is convex, so $V$ is contractible.
It follows that $U$ is acyclic.
\end{proof}

\end{document}
